% Architectural projection of WO-053-A.
% This document intentionally omits implementation ordering and tracker state.
\documentclass[11pt]{article}

\usepackage[T1]{fontenc}
\usepackage[utf8]{inputenc}
\usepackage{lmodern}
\usepackage{microtype}
\usepackage[margin=1in]{geometry}
\usepackage{enumitem}
\usepackage{xcolor}
\usepackage[hidelinks]{hyperref}

\definecolor{leafgreen}{HTML}{2F6B45}
\definecolor{leafgray}{HTML}{4B5563}
\newcommand{\LeafOS}{\textsc{LeafOS}}
\newcommand{\contract}[1]{\texttt{#1}}

\title{\textbf{\LeafOS{} Model Pack and Installer Trust Architecture}\\
\large A design mission for reliable local-model acquisition and discovery}
\author{LeafOS Architecture Series}
\date{August 9, 2026}

\begin{document}

\maketitle

\begin{abstract}
The local-model installer is the first operational trust boundary between a
declared model architecture and the weights that later enter execution. Its
purpose is not merely to move files. It must turn model intent into a stable,
inspectable, and verifiable local state that every LeafOS surface understands
in exactly the same way. This document defines that mission, the architectural
boundaries that support it, and the truth conditions required before the
llama.cpp capability preflight may consume a model path.
\end{abstract}

\section{Mission Statement}

The \LeafOS{} model-pack and installation system shall provide one coherent
answer to five questions:

\begin{quote}
\color{leafgreen}\bfseries
What model was selected?  Why was it selected?  Where will it live?  What
operation is permitted next?  What evidence proves the answer?
\end{quote}

The guided interface, canonical backend, model catalog, pack manifests, local
store, runtime status, and capability preflight must agree on those answers.
Convenience surfaces may explain or visualize the process, but they may not
reinterpret success, hide failure, invent a destination, or substitute a stale
catalog for the current source of truth.

The mission is therefore to make model acquisition \emph{deterministic in
contract, explicit in authority, resumable in operation, and honest in
evidence}. A user should be able to inspect the system before permitting a
large transfer, interrupt it without losing durable progress, and return later
to the same model identity and destination without reconstructing intent from
conversation or shell history.

\section{Architectural Premise}

Model installation is a state transition, not a download shortcut. The
canonical lifecycle is:

\begin{verbatim}
catalog -> offline plan -> metadata resolution -> explicit apply -> verification
\end{verbatim}

Each boundary has a distinct meaning.

\begin{description}[style=nextline,leftmargin=0pt,labelindent=0pt]
  \item[Catalog]
  Defines the allowable repositories, model identities, roles, variants,
  quantizations, expected artifact patterns, and safety classification.

  \item[Offline plan]
  Expands user intent into a reviewable set of model and quantization choices,
  a deterministic destination, and an estimated resource commitment. Planning
  must not require a network connection or begin a transfer.

  \item[Metadata resolution]
  Reads remote metadata, selects exact files, and pins immutable repository
  revisions. It changes knowledge, not local model weights.

  \item[Explicit apply]
  Crosses the transfer boundary only after acknowledgement. It owns download,
  resume, locking, bounded concurrency, and durable event recording.

  \item[Verification]
  Establishes whether the local artifacts satisfy the resolved contract. It
  must be distinct from inventory and must never infer validity from presence
  alone.
\end{description}

The lifecycle is shared by standard profiles and model packs. A pack adds role
composition and orchestration intent, but it does not create a second download
engine, a second catalog, or a weaker acceptance path.

\section{One Truth Across Surfaces}

The Python installation backend is the semantic authority for planning,
resolution, application, and verification. PowerShell, Bash, menus, launchers,
and dashboards are projections over that authority. Their responsibilities are
limited to argument transport, presentation, confirmation, and faithful exit
propagation.

This yields four invariants:

\begin{itemize}[leftmargin=*]
  \item Arguments reach the backend as exact tokens, including paths containing
  spaces and non-default storage locations.
  \item A backend failure remains a failure at every wrapper boundary, with the
  original reason preserved for both humans and automation.
  \item Installed tooling and source tooling identify the same catalog version,
  command set, profiles, slots, and catalog digest before readiness is granted.
  \item Human-readable and structured output describe the same state and legal
  next action.
\end{itemize}

Readiness is consequently a parity claim, not merely evidence that Python is
installed and a directory is writable. A stale environment must refuse to
present itself as ready and must provide one deterministic repair action.

\section{Model Packs as Typed Intent}

A model pack is a typed declaration of roles and the catalog artifacts needed
to realize them. The installer must distinguish among installable packs,
identity-only architectural packs, companion manifests, and invalid documents.
File extension alone is not a sufficient type system.

An installable pack must name a supported schema and contain a non-empty set of
catalog-backed installation items. Each item resolves through the canonical
catalog. Multiple roles may share one model and quantization; the plan may
deduplicate the physical transfer while retaining the role-level intent needed
by the runtime.

Pack identity does not weaken catalog policy. A pack cannot introduce an
unregistered repository, silently broaden an artifact pattern, or relabel an
experimental model as ordinary. Unsupported documents remain inspectable as
design material but are not offered as executable installation choices.

\section{Deterministic Model-Root Discovery}

Every surface must resolve the local model root through one precedence
contract. The selected root is an architectural fact and must be reported with
its reason. The contract recognizes:

\begin{itemize}[leftmargin=*]
  \item an explicit destination supplied for the current operation;
  \item the \contract{LEAF\_MODEL\_DIR} environment contract;
  \item a recognized installer-local model store when present; and
  \item the documented per-user default as the final fallback.
\end{itemize}

The same result must appear in doctor output, plans, pack plans, status,
verification, runtime provider configuration, and the handoff to WO-053. No
surface may create a second model library merely because another surface failed
to expose the already populated one.

The resolved root should be absolute and accompanied by a machine-readable
reason. Plans bind this destination so later resolution and application cannot
quietly drift to another volume or user profile.

\section{Safety and Experimental Authority}

Large model transfers are consequential operations. Safety must be represented
in the plan rather than added as a final prompt. Estimates, free-space margins,
artifact counts, fallback decisions, and experimental status must be visible
before application.

Base packs and experimental extensions are separate authorities. Dormant or
experimental lanes must not enter a base plan implicitly. When an operator
chooses an experimental plan, every experimental catalog key in that plan must
be acknowledged; confirmation of one key cannot authorize unrelated heavy
models by association.

This policy preserves useful experimentation without making the ordinary path
ambiguous. Experimental models remain first-class catalog entries with explicit
identity and evidence, but their cost and risk cannot be hidden behind role
deduplication or friendly pack naming.

\section{Verification Doctrine}

Inventory answers \emph{what files are present}. Verification answers
\emph{whether those files satisfy a declared contract}. The two operations must
have different names, fields, and exit behavior.

A verified GGUF store requires, as applicable:

\begin{itemize}[leftmargin=*]
  \item valid GGUF magic and a credible minimum artifact size;
  \item catalog membership and an expected model/variant relationship;
  \item exact resolved paths and byte sizes when a pinned plan is available;
  \item explicit treatment of missing, duplicate, unknown, or unclassified
  artifacts; and
  \item a truthful aggregate disposition derived from every required artifact.
\end{itemize}

Optional cryptographic hashes strengthen identity and provenance, but a hash is
not a substitute for classification. Conversely, catalog recognition is not a
substitute for local integrity. Verification combines structural validity,
resolved expectation, and evidence provenance.

The verifier fails closed. Zero-byte placeholders, invalid headers, missing
pinned files, size mismatches, and unexplained artifacts cannot produce a green
result merely because the directory exists.

\section{Relationship to the llama.cpp Capability Boundary}

This architecture prepares reliable model identity and location. It does not
prove that llama.cpp can load the selected model, that the requested CPU or GPU
backend is available, or that inference produces tokens.

WO-053 consumes the verified handoff and independently proves executable
identity, GGUF compatibility, backend devices, host facts, and capability
policy. WO-054 through WO-056 remain responsible for live generation,
streaming, cancellation, timeout, cleanup, and execution evidence.

The boundary is deliberate:

\begin{verbatim}
installer verification  ->  artifact is the declared local artifact
capability preflight     ->  artifact and backend may enter evaluation
live inference           ->  fresh tokens were generated by owned execution
\end{verbatim}

None of these claims implies the next. Each produces evidence appropriate to
its own authority.

\section{Design Success}

The architecture succeeds when a new or returning user can inspect a pack,
understand its resource and safety implications, produce an offline plan,
resolve exact immutable artifacts, apply or resume deliberately, verify the
result, and hand one unambiguous model path into capability evaluation.

Success is observable when wrappers and backend agree, stale tooling refuses
readiness, unsupported pack documents are classified rather than executed,
storage discovery is consistent, experimental authority is explicit, and
corrupt fixtures cannot appear verified. These properties must be demonstrated
with temporary environments and synthetic artifacts so validation does not
depend on downloading or hashing the production model library.

\section{Non-Goals}

This design does not choose model quality rankings, alter runtime role policy,
assign GPU placement, perform remote transfers as part of validation, or start
an inference provider. It does not claim that a structurally valid model is
semantically good, operationally compatible, or safe for an unrestricted
workload. Those decisions belong to later capability and evaluation gates.

\section{Conclusion}

The model installer is the bridge between architectural intent and executable
local state. If that bridge offers multiple catalogs, ambiguous destinations,
decorative verification, or softened failure codes, every downstream claim is
weaker. If it is deterministic, typed, explicit, and evidence-bearing, WO-053
can begin from a trustworthy artifact rather than a guessed path.

The governing principle is simple:

\begin{quote}
\color{leafgreen}\bfseries
The installer does not prove intelligence. It proves that the exact model we
intended is the model we possess, in the place every LeafOS component agrees to
find it.
\end{quote}

\vfill
\begin{center}
\small\color{leafgray}
Architectural projection derived from WO-053-A. Implementation ordering,
tracker state, and evidence chronology remain in the work-order record.
\end{center}

\end{document}
